\chapter*{Introduzione al Corso}

Quantum information è un corso sulla teoria dell'informazione che si compone di 2 parti
\begin{itemize}
	\item 3 CFU di Introduzione alla \textbf{Teoria dell'Informazione Classica}
	\item 3 CFU di Introduzione alla \textbf{Teoria dell'Informazione Quantistica}
\end{itemize}
Cos'è la Teoria dell'Informazione?\\
Fino a che punto possiamo comprimere i dati?\\
Qual'è il miglior tasso di trasmissione possibile?
\begin{itemize}
	\item \textbf{Teoria classica dell'informazione}
	\\ Ci chiediamo cosa sia l'informazione e come renderla affidabile su canali inaffidabili. Vediamo in breve cenni di teoria dell'informazione di Shannon:
	\begin{itemize}
		\item È possibile inviare informazione con probabilità di errore $\varepsilon$, se il tasso di comunicazione è minore della capacità del canale
		\item Ogni segnale ha complessità intrinseca (\textit{entropia}) oltre al quale non può essere compreso
		\item Se l'entropia della sorgente è minore della capacità di canale, una comunicazione (asintoticamente) priva di errori può avvenire
	\end{itemize}
	\item  \textbf{Teoria dell'informazione quantistica}
	\\A differenza del caso precedente utilizziamo i qubit. Quello che fa davvero la differenza sono i \textit{termini di coerenza}, usando le matrici di densità. 
	Un'altra differenza è la presenza dell'\textbf{entanglement}.
	Studieremo come misurare l'entropia di uno stato quantistico (Entropia di Von Neumann).
\end{itemize}
Vedremo le seguenti misure: Capacità classica, quantistica e quantistica assistita (con entanglement).
